% =====================================================================
%  2 UŽDUOTIS — DI ir mašininio mokymosi metodų taikymo galimybės
%  Rugsėjo 2 d. | Tikslinė apimtis: 5–6 psl. (RIBA 6,5)
%
%  REIKALINGA PREAMBULĖJE: \usepackage{xltabular}
%
%  Būklė: tab:metodai užpildyta (T3). Poskyrių tekstas — T2 medžiaga
%  faile rezultatai/darbiniai/metodu_apzvalga.md
% =====================================================================

\subsection{Vertinimo rėmas ir apžvalgos ašis}

Pirmame skyriuje iš IoT aplinkos savybių išvesti reikalavimai aptikimo
sprendimui (žr. \ref{tab:reikalavimai} lentelę). Šiame skyriuje jie
nekuriami iš naujo --- tik pritaikomi konkretiems metodams. Toks
paveldėjimas svarbus metodologiškai: kriterijai suformuluoti dar
nežinant, kurie metodai bus svarstomi, todėl atranka negali būti
pritempta prie iš anksto pageidaujamo atsakymo.

Metodai grupuojami pagal mokymosi paradigmą, o ne pagal algoritmų
šeimą. Tokį pasirinkimą lemia tai, kad metodo tinkamumą šiame darbe
labiausiai apibrėžia ne jo matematinė prigimtis, o tai, ko jis
reikalauja iš duomenų: pažymėtų pavyzdžių, sekos, tinklo topologijos ar
paskirstytų klientų. Tolesniuose poskyriuose matyti, kad būtent šis
reikalavimas dažniausiai ir nulemia, ar metodą apskritai galima
taikyti turimam duomenų rinkiniui.

\subsection{Prižiūrimas mokymasis}

Prižiūrimi metodai reikalauja pažymėtų mokymo pavyzdžių. Naudojamame
rinkinyje jų yra keliasdešimt milijonų, todėl paradigma taikytina be
išlygų, o jos riba yra kita: modelis atpažįsta tik tas atakas, kurių
pavyzdžių matė mokydamasis.

Lentelinei srauto statistikai atskaitos klasė yra sprendimų medžių
ansambliai. Jie skiriasi tuo, kaip derinami atskiri medžiai:
atsitiktinis miškas juos moko lygiagrečiai iš skirtingų imties poaibių
ir apjungia balsavimu, o gradientinio stiprinimo modeliai (XGBoost,
LightGBM) mokomi nuosekliai, kiekvienam taisant ankstesniųjų liekamąją
paklaidą. Abi kryptys tvarko klasių disbalansą svoriais, nedubliuodamos
duomenų, ir abi inferuoja per mikrosekundes
\cite{almahaqeri2026gradient}.

Kiti prižiūrimi metodai IoT kontekste susiduria su tuo pačiu
sunkumu --- jų kaina atsiduria ten, kur jos leisti negalima. Atraminių
vektorių mašinos mokymo sudėtingumas auga kvadratiškai ar kubiškai
imties atžvilgiu, o artimiausių kaimynų metodas mokymo fazės neturi,
tačiau kiekvienam sprendimui turi peržiūrėti visą mokymo aibę --- taigi
lėčiausias yra būtent tada, kai reikia atsakyti realiu laiku. Naivųjį
Bajeso klasifikatorių, pigiausią iš visų, atmeta jo paties prielaida:
požymių sąlyginis nepriklausomumas šiame rinkinyje paneigiamas
tiksliomis funkcinėmis priklausomybėmis tarp statistinių požymių.

\subsection{Neprižiūrimas mokymasis ir anomalijų aptikimas}

Neprižiūrimi metodai mokosi tik iš gerybinio srauto ir atakos pavyzdžių
nereikalauja. Tai vienintelė paradigma, iš principo galinti aptikti
nematytas grėsmes, todėl darbe, kurio objektas yra prevencija, ji
negali būti praleista. Naudojamame rinkinyje gerybinio srauto yra
daugiau nei milijonas įrašų --- mokymui su kaupu pakanka.

Autokoderis suspaudžia įvestį per siaurą sluoksnį ir mokosi ją atkurti;
neįprastas srautas atkuriamas prasčiau, todėl atkūrimo paklaida veikia
kaip anomalijos matas. Toks sprendimas realiame IoT sraute jau
demonstruotas visoje grandinėje nuo požymių išgavimo iki aptikimo
\cite{meidan2018nbaiot}. Izoliacijos miškas tą pačią užduotį sprendžia
kur kas pigiau: atsitiktiniai skaidymai anomalijas atskiria greičiau nei
tankiai išsidėsčiusius įprastus taškus. Vienos klasės atraminių vektorių
mašina koncepciškai artima, bet paveldi tą patį mokymo sudėtingumą kaip
ir prižiūrima jos atmaina. Klasterizavimo metodai duoda grupes, o ne
sprendimą, todėl tiesiogiai palyginti su kitais metodais jų nepavyktų.

Šiai paradigmai klasių disbalansas nėra taikytinas kriterijus. Modelis
mato tik gerybinį srautą, todėl santykiai tarp atakų klasių jo
mokymo nepasiekia --- tai ne pranašumas, o kitaip suformuluotas
uždavinys.

\subsection{Gilusis mokymasis}

Gilaus mokymosi patrauklumas grindžiamas gebėjimu pačiam išmokti
požymių reprezentaciją. Tačiau kiekviena architektūra tai daro
remdamasi prielaida apie duomenų sandarą, ir būtent čia paaiškėja
esminis šio darbo apribojimas: naudojamo rinkinio įrašas yra trisdešimt
šešios agreguotos statistikos, be eilės, be krypties ir be topologijos.

Konvoliuciniai tinklai remiasi prielaida, kad gretimos įvesties
reikšmės susijusios, tačiau požymių eilė lentelėje yra sutartinė.
Rekurentiniai tinklai ir dėmesio mechanizmu grįstos architektūros
modeliuoja seką, o rinkinyje nėra nei laiko žymos, nei eilės tarp
įrašų --- kiekvienas įrašas jau yra dešimties arba šimto paketų lango
apibendrinimas, taigi laiko matmuo į požymius jau suvyniotas. Grafų
neuroniniams tinklams reikia mazgų ir briaunų, o adresų ar krypties
rinkinyje nėra. Vienintelė architektūra, nedaranti nė vienos iš šių
prielaidų, yra daugiasluoksnis perceptronas; ribotų išteklių
įrenginiuose jis kartu su laikine konvoliucija rodo geriausią tikslumo
ir kainos santykį \cite{mazinani2026constrained}.

Iš to seka klausimas, kurį šis darbas ir tikrina: ar didesnis modelio
sudėtingumas duoda naudos, kai įvestis yra nedidelis skaičius
nepriklausomų statistikų.

\subsection{Paskirstytas mokymasis ir paaiškinamumas}

Federuotas mokymasis nėra atskiras algoritmas, o mokymo schema: modeliai
mokomi lokaliai, o centralizuojami tik jų svoriai, todėl srautas
nepalieka tinklo. IoT kontekste tai reikšminga privatumo požiūriu, o
ryšio kaina gali būti optimizuota \cite{nassef2026tinyml}. Schemai
reikia duomenų, padalytų į klientus pagal realų šaltinį; naudojamame
rinkinyje įrenginio identifikatoriaus nėra, todėl atsitiktinis
padalijimas imituotų vienodą pasiskirstymą ir tirtų ne tą reiškinį,
dėl kurio schema ir taikoma.

Paaiškinamumas nėra mokymosi paradigma, o skersinis sluoksnis. SHAP
metodas pateikia nuoseklias globalias ir lokalias požymių atributikas,
LIME --- greitesnius lokalius paaiškinimus pavieniams pranešimams
audituoti; abu reikalauja pastebimų skaičiavimo išteklių ir todėl
netinka realaus laiko grandinei \cite{mohale2025xai}. Dėl to šiame
darbe paaiškinamumas naudojamas kaip rezultatų analizės priemonė, o ne
kaip aptikimo sprendimo dalis. Ši skirtis atsispindi \ref{tab:metodai}
lentelės stulpelyje: savaime interpretuojami metodai sprendimą leidžia
parodyti be papildomos kainos, o ansambliams ir neuroniniams tinklams ta
pati galimybė pasiekiama tik atskiru skaičiavimu.


% =====================================================================
%  tab:metodai — metodų vertinimo matrica (T3)
% ---------------------------------------------------------------------
%  Kaip ir tab:atakos: xltabular NEGALI būti table float'e,
%  todėl \caption rašomas lentelės viduje.
%
%  Kiekvienas įvertinimas turi šaltinį arba žymą (aut.) — tuščių
%  langelių stulpelyje „Šalt.“ būti negali. Žr. plano 4 skyrių.
% =====================================================================

\begingroup
\footnotesize
\setlength{\tabcolsep}{3pt}
\sloppy
\begin{xltabular}{\textwidth}{@{}
  >{\raggedright\arraybackslash}p{2.8cm}
  >{\raggedright\arraybackslash}p{1.35cm}
  >{\centering\arraybackslash}p{1.0cm}
  >{\centering\arraybackslash}p{0.95cm}
  >{\centering\arraybackslash}p{1.05cm}
  >{\raggedright\arraybackslash}p{1.55cm}
  >{\raggedright\arraybackslash}X
  >{\raggedright\arraybackslash}p{1.7cm}@{}}
\caption{Dirbtinio intelekto metodų vertinimas pagal reikalavimus,
         kylančius iš IoT aplinkos savybių}
\label{tab:metodai} \\
\toprule
\textbf{Metodas} & \textbf{Para\-digma} & \textbf{Nemat. atakos}
  & \textbf{Moky\-mo kaina} & \textbf{Infe\-ren. kaina}
  & \textbf{Interpre\-tuojamumas} & \textbf{Atsparumas disbalansui}
  & \textbf{Šalt.} \\
\midrule
\endfirsthead
\multicolumn{8}{@{}l}{\itshape \ref{tab:metodai} lentelės tęsinys} \\
\toprule
\textbf{Metodas} & \textbf{Para\-digma} & \textbf{Nemat. atakos}
  & \textbf{Moky\-mo kaina} & \textbf{Infe\-ren. kaina}
  & \textbf{Interpre\-tuojamumas} & \textbf{Atsparumas disbalansui}
  & \textbf{Šalt.} \\
\midrule
\endhead
\bottomrule
\endlastfoot

% --- Prižiūrimas mokymasis ---
Sprendimų medis (CART) & Prižiūr. & Ne & maž. & maž. &
  savaiminis & silpnas; taisomas svoriais & (aut.) \\
Random Forest & Prižiūr. & Ne & vid. & maž. &
  per XAI & geras su \texttt{class\_weight} & \cite{imani2025imbalance} \\
XGBoost & Prižiūr. & Ne & vid. & maž. &
  per XAI & geras; \SI{1.05}{\micro\second} inferencija &
  \cite{almahaqeri2026gradient} \\
LightGBM & Prižiūr. & Ne & maž. & maž. &
  per XAI & geras; beveik lygu XGBoost & \cite{almahaqeri2026gradient} \\
SVM (RBF) & Prižiūr. & Ne & l. did. & vid. &
  juod. dėžė & vid.; $O(n^2)$--$O(n^3)$ mokymas & (aut.) \\
$k$ artimiausių kaimynų & Prižiūr. & Ne & --- & l. did. &
  savaiminis & prastas; kaimynystėje dominuoja gausios klasės &
  \cite{sallam2026gap} \\
Naive Bayes & Prižiūr. & Ne & maž. & maž. &
  savaiminis & vid.; prielaida pažeista & (aut.) \\
\addlinespace
% --- Neprižiūrimas mokymasis ---
Autokoderis & Neprižiūr. & \textbf{Taip} & vid. & maž. &
  juod. dėžė & neaktualus --- mokoma tik iš gerybinio srauto &
  \cite{meidan2018nbaiot} \\
Isolation Forest & Neprižiūr. & \textbf{Taip} & maž. & maž. &
  per XAI & neaktualus --- kaip ir autokoderio & (aut.) \\
One-Class SVM & Neprižiūr. & \textbf{Taip} & l. did. & vid. &
  juod. dėžė & neaktualus, bet mokymas netelpa & (aut.) \\
Klasterizavimas ($k$-vid., DBSCAN) & Neprižiūr. & Ribotai & vid. & vid. &
  per XAI & prastas; mažos grupės susilieja & (aut.) \\
\addlinespace
% --- Pusiau prižiūrimas ---
Savimoka (pseudo-žymėjimas) & Pusiau prižiūr. & Ne & vid. & maž. &
  kaip bazinio & \textbf{blogėja} --- pseudo-etiketės stiprina
  gausių klasių šališkumą & (aut.) \\
\addlinespace
% --- Gilusis mokymasis ---
Daugiasluoksnis perceptronas & Gilusis & Ne & vid. & maž. &
  juod. dėžė & vid.; svoriai nuostolio funkcijoje &
  \cite{mazinani2026constrained} \\
1D-CNN & Gilusis & Ne & did. & vid. &
  juod. dėžė & vid.; požymių eilė savavališka &
  \cite{mazinani2026constrained} \\
LSTM / GRU & Gilusis & Ne & did. & vid. &
  juod. dėžė & vid.; sekos duomenyse nėra &
  \cite{mazinani2026constrained} \\
Transformer & Gilusis & Ne & l. did. & did. &
  juod. dėžė & vid.; ta pati sekos problema & \cite{nassef2026tinyml} \\
Grafų neuroninis tinklas & Gilusis & Ribotai & l. did. & did. &
  juod. dėžė & vid.; \SIrange{120}{180}{\milli\second} Raspberry Pi 4 &
  \cite{nassef2026tinyml} \\
\addlinespace
% --- Paskirstytas mokymasis ---
Federuotas mokymasis (FedAvg)\textsuperscript{a} & Paskirst. & kaip bazinio &
  paskirst. & kaip bazinio & kaip bazinio &
  \textbf{blogėja} esant ne-IID duomenims & \cite{nassef2026tinyml} \\
\end{xltabular}

\vspace{-0.5em}
{\footnotesize
\textbf{Žymėjimai.} Kaina: \emph{maž.} — mažas, \emph{vid.} — vidutinis,
\emph{did.} — didelis, \emph{l.\ did.} — labai didelis;
„---“ reiškia, kad fazės nėra.
Interpretuojamumas: \emph{savaiminis} — sprendimą galima parodyti be
papildomų įrankių; \emph{per XAI} — reikia SHAP ar LIME, o tai
papildoma skaičiavimo kaina \cite{mohale2025xai};
\emph{juod.\ dėžė} — savaiminio paaiškinimo nėra.
Žyma \emph{(aut.)} nurodo, kad įvertinimas yra šio darbo autoriaus, o ne
perimtas iš šaltinio.
\textsuperscript{a}\,Federuotas mokymasis yra mokymo schema, o ne atskiras
algoritmas: jo savybės priklauso nuo bazinio modelio.
\par}
\endgroup

% ---------------------------------------------------------------------
% 2.6-2.9   (T5, T4, T6, T7)
% ---------------------------------------------------------------------
\subsection{Praktiniai apribojimai IoT aplinkoje}
\label{sec:apribojimai}

Ankstesniuose poskyriuose metodai vertinti pagal aptikimo galimybes.
Šis poskyris prideda antrą matmenį --- ar metodas apskritai gali veikti
ten, kur jį reikėtų įdiegti. Trys apribojimai yra kiekybiniai ir todėl
tinka kaip atrankos kriterijai: inferencijos delsa, mokymo kaina ir
klasių disbalansas.

\subsubsection{Inferencijos delsa ir diegimo vieta}

Pirmame skyriuje diegimo vieta pasirinkta kraštiniame šliuze, o jam
keliamas realaus laiko reikalavimas yra \SIrange{20}{50}{\milli\second}
vienam sprendimui; įrenginio lygmeniui skiriama
\SIrange{1}{10}{\milli\second}, debesiui ---
ne mažiau kaip \SI{100}{\milli\second} \cite{sallam2026gap}. Šie skaičiai leidžia
metodus vertinti ne abstrakčiai, o pagal tai, kiek biudžeto jie suvartoja.

Skirtumas tarp lentelinių ansamblių ir giliųjų architektūrų yra ne
laipsnio, o eilės. Gradientinio stiprinimo modelis tame pačiame
CICIoT2023 rinkinyje 34 klasių užduotyje inferuoja per
\SI{3.53}{\micro\second} vienam įrašui \cite{almahaqeri2026gradient} ---
tai \num{0.018}~\% griežtesniojo (\SI{20}{\milli\second}) biudžeto.
Hibridinis grafų ir rekurentinis modelis tame pačiame rinkinyje
reikalauja \SIrange{120}{180}{\milli\second} Raspberry~Pi~4 plokštėje
\cite{nassef2026tinyml}, t.\,y. nuo 2,4 iki 9 kartų \emph{daugiau} nei
leidžia biudžetas. Tarp šių dviejų verčių yra keturios eilės.

\begin{table}[htbp]
\centering
\caption{Metodų tinkamumas pagal diegimo vietos delsos biudžetą}
\label{tab:apribojimai}
\footnotesize
\setlength{\tabcolsep}{4pt}
\begin{tabularx}{\textwidth}{@{}
  >{\raggedright\arraybackslash}p{2.6cm}
  >{\raggedright\arraybackslash}p{2.0cm}
  >{\raggedright\arraybackslash}X
  >{\raggedright\arraybackslash}X@{}}
\toprule
\textbf{Diegimo vieta} & \textbf{Delsos biudžetas} & \textbf{Telpa}
  & \textbf{Netelpa} \\
\midrule
Įrenginys (mikro\-valdiklis) & \SIrange{1}{10}{\milli\second} &
  Kvantuotas MLP arba TCN (INT8, modelis mažesnis
  \textgreater\,90~\%) \cite{mazinani2026constrained} &
  Ansambliai --- ne dėl delsos, o dėl modelio dydžio atmintyje \\
\addlinespace
\textbf{Kraštinis šliuzas} & \SIrange{20}{50}{\milli\second} &
  XGBoost, LightGBM, Random Forest, Isolation Forest, MLP ---
  su kelių eilių atsarga &
  Hibridiniai gilieji modeliai
  (\SIrange{120}{180}{\milli\second}) \cite{nassef2026tinyml} \\
\addlinespace
Debesis & $\geq$~\SI{100}{\milli\second} & Visi apžvelgti metodai &
  Prarandamas reagavimas realiu laiku; srautas turi palikti tinklą \\
\bottomrule
\end{tabularx}
\end{table}

Skelbiamos vertės išmatuotos skirtinga aparatūra: mikrosekundinės ---
serverio klasės procesoriumi, \SIrange{120}{180}{\milli\second} ---
vienos plokštės kompiuteriu. Skirtumo eilė tai atlaiko:
\SI{3.53}{\micro\second} į \SI{20}{\milli\second} biudžetą telpa
\num{5669} kartus, todėl reikalavimą tenkintų ir tiek kartų lėtesnė
aparatūra.

Atskirai pažymėtina, kad neprižiūrimų metodų atveju skelbiamas dydis
dažnai yra ne inferencijos delsa, o \emph{aptikimo laikas}. Gilusis
autokoderis N-BaIoT rinkinyje ataką aptinka per
$174 \pm \SI{212}{\milli\second}$ \cite{meidan2018nbaiot}: šis dydis
apima ir srauto lango sukaupimą, ne vien modelio skaičiavimą.
Standartinis nuokrypis čia didesnis už vidurkį, todėl vien vidurkiu
remtis negalima. Todėl penktame skyriuje matuojama grynoji inferencijos delsa, o lango
sukaupimo laikas nurodomas atskirai.

\subsubsection{Mokymo kaina}

Mokymo kaina yra ne vienkartinis, o pasikartojantis dydis. Tinklo
elgsena kinta, todėl aptikimo modelis turi būti periodiškai
permokomas, o kraštinio šliuzo aplinkoje tam skiriama įprasta
aparatūra be grafinio spartintuvo. Dėl to mokymo sudėtingumas yra
atrankos kriterijus, o ne vien techninė aplinkybė.

Pirmiausia jis pašalina metodus, kurių mokymo sudėtingumas imties
atžvilgiu yra kvadratinis ar kubinis, --- atraminių vektorių mašiną ir
jos vienos klasės atmainą. Milijonų eilučių eilės imtyse toks
augimas mokymą paverčia nepraktišku nepriklausomai nuo turimos
aparatūros. Giliosioms architektūroms mokymo trukmė be spartintuvo
veikia kaip antras, nepriklausomas argumentas prie jau įvardyto
delsos perviršio.

\subsubsection{Klasių disbalansas}

Naudojamame rinkinyje disbalansas reiškiasi dviem skirtingais
santykiais, ir jie turi skirtingas pasekmes. Atakų ir gerybinio srauto
santykis yra $41{,}8:1$ --- dėl jo bendras tikslumas dvejetainėje
užduotyje tampa neinformatyvus, nes viską ataka žymintis modelis
pasiektų apie $97{,}7$~\%. Kur kas didesnis yra santykis tarp pačių
klasių: gausiausia klasė už rečiausią didesnė \num{5764} kartus.
Būtent antrasis santykis paaiškina, kodėl daugiaklasėje užduotyje
krinta makro vidurkiu skaičiuojamas F1 --- septynios klasės sudaro
mažiau nei $0{,}03$~\% duomenų ir kiekviena jų į makro vidurkį įeina
tokiu pat svoriu kaip milijoninės.

Tai patvirtina ir literatūra. Tame pačiame CICIoT2023 rinkinyje
gradientinio stiprinimo modelio tikslumas pereinant nuo dvejetainės
prie aštuonių klasių užduoties beveik nekinta ($99{,}61 \rightarrow
99{,}59$~\%), o makro-F1 krinta nuo $0{,}9952$ iki $0{,}8903$
\cite{almahaqeri2026gradient}. Tas pats modelis ir tie patys duomenys ---
skiriasi tik metrika ir užduoties detalumas.

Balansavimo priemonių palyginimas prie skirtingų disbalanso lygių rodo,
kad suderintas gradientinio stiprinimo modelis kartu su SMOTE metodu
duoda nuosekliai geriausią F1, o atsitiktinio miško tikslumas esant
stipriam disbalansui pastebimai krinta \cite{imani2025imbalance}. Tas
pats šaltinis pateikia ir metodinę pastabą, svarbią penktam skyriui:
didėjant disbalansui F1, Kappa ir Matthewso koreliacijos koeficientas
svyruoja stipriai, o ROC-AUC ir PR-AUC išlieka palyginti stabilūs.
Todėl šalia makro-F1 tikslinga pateikti ir PR-AUC.

\subsection{Susiję darbai ir skelbiami rezultatai}

Skelbiamų rezultatų apžvalga reikalinga ne tam, kad būtų surinkti aukšti
skaičiai, o tam, kad savo matavimams atsirastų atskaitos taškas.
\ref{tab:susije} lentelėje kiekvienam darbui nurodoma ne tik metrika,
bet ir užduoties detalumas bei tai, kas riboja rezultato aiškinimą.

Užduoties detalumo stulpelis yra būtinas, nes be jo skaičiai
nepalyginami. Tai geriausiai matyti to paties modelio eilutėse: tame
pačiame duomenų rinkinyje gradientinio stiprinimo modelio tikslumas
dvejetainėje ir aštuonių klasių užduotyse skiriasi vos šimtosiomis
procento dalimis, o makro vidurkiu skaičiuojamas F1 krinta nuo
$0{,}9952$ iki $0{,}8903$ \cite{almahaqeri2026gradient}. Darbas,
pateikiantis tik tikslumą ir nenurodantis klasių skaičiaus, tokio
skirtumo neatskleidžia.

Antra pastebima savybė --- ne visos abejonės kyla dėl silpnų metrikų.
Dalis darbų skelbia aukštus F1 įverčius, tačiau užduoties detalumo
nenurodo \cite{nassef2026tinyml}, kitas pateikia suminį kompromiso
įvertį vietoj metrikų kiekvienam modeliui \cite{mazinani2026constrained},
o klasikiniame botnetų aptikimo darbe pasiektas šimtaprocentinis
teisingai atpažintų atakų dažnis atitinka lengviausią įmanomą
formuluotę --- atskiras modelis kiekvienam įrenginiui ir dvi botnetų
šeimos \cite{meidan2018nbaiot}. Trūkstamas metodikos elementas riboja
palyginimą lygiai taip pat, kaip ir prastas rezultatas.

Sugretinant savo rezultatus su šiais darbais svarbu turėti omenyje
požymių aibės skirtumą. Dauguma publikacijų remiasi pilna keturiasdešimt
šešių požymių versija, o šiame darbe naudojamame leidime jų
trisdešimt devyni: pašalinti krypties požymiai ir iš jų išvestos
statistikos. Dalis literatūroje aukštai reitinguojamų požymių šiame
leidime neegzistuoja, todėl skelbiami įverčiai laikytini bendru
atskaitos lygiu, o ne tiesioginiu palyginimo tašku. Į tai atsižvelgiama
šeštame skyriuje.

Iš apžvalgos taip pat matyti, kokio dydžio rezultato realu tikėtis.
Aštuonių kategorijų užduotyje geriausias literatūroje skelbiamas makro-F1
yra apie $0{,}89$; tai ir yra vertė, prie kurios lyginami šio darbo
matavimai --- ne artimi vienetui tikslumo įverčiai, gaunami dvejetainėje
formuluotėje.

% =====================================================================
%  tab:susije -- susije darbai ir skelbiami rezultatai (T4)
% ---------------------------------------------------------------------
%  Stulpelis "Uzduotis" yra butinas: be jo metrikos nepalyginamos.
%  Stulpelis "Kas kelia abejoniu" yra visa lenteles verte -- be jo
%  tai citatu sarasas, ne analize.
%
%  BUKLE 2026-09-02: 8 eilutes is 3 saltiniu -- visi skaiciai patikrinti
%  leidejo puslapyje. TRUKSTA 4 eiluciu (zr. komentara po lentele).
% =====================================================================

\begingroup
\footnotesize
\setlength{\tabcolsep}{3pt}
\sloppy
\begin{xltabular}{\textwidth}{@{}
  >{\raggedright\arraybackslash}p{2.0cm}
  >{\raggedright\arraybackslash}p{2.0cm}
  >{\raggedright\arraybackslash}p{2.1cm}
  >{\raggedright\arraybackslash}p{1.5cm}
  >{\raggedright\arraybackslash}p{2.8cm}
  >{\raggedright\arraybackslash}X@{}}
\caption{Skelbiami dirbtinio intelekto metodų rezultatai IoT įsibrovimų
         aptikimo darbuose}
\label{tab:susije} \\
\toprule
\textbf{Šaltinis} & \textbf{Duomenų rinkinys} & \textbf{Metodas}
  & \textbf{Užduotis} & \textbf{Skelbiama metrika}
  & \textbf{Kas kelia abejonių} \\
\midrule
\endfirsthead
\multicolumn{6}{@{}l}{\itshape \ref{tab:susije} lentelės tęsinys} \\
\toprule
\textbf{Šaltinis} & \textbf{Duomenų rinkinys} & \textbf{Metodas}
  & \textbf{Užduotis} & \textbf{Skelbiama metrika}
  & \textbf{Kas kelia abejonių} \\
\midrule
\endhead
\bottomrule
\endlastfoot

\cite{almahaqeri2026gradient} & CICIoT2023 (46 pož.) & XGBoost &
  dvejetainė & tikslumas 99,61~\%; macro-F1 0,9952;
  \SI{0.36}{\micro\second}/įrašui &
  Naudota pilna 46 požymių aibė; šio darbo leidime jų 39 \\
\cite{almahaqeri2026gradient} & CICIoT2023 (46 pož.) & XGBoost &
  8 klasės & tikslumas 99,59~\%; \textbf{macro-F1 0,8903};
  \SI{1.05}{\micro\second} &
  Tikslumas beveik nepakito, o macro-F1 krito 0,10 --- tas pats modelis,
  tie patys duomenys \\
\cite{almahaqeri2026gradient} & CICIoT2023 (46 pož.) & XGBoost &
  34 klasės & tikslumas 99,48~\%; macro-F1 0,8876;
  \SI{3.53}{\micro\second} &
  Tas pats dėsningumas ryškėja dar labiau \\
\cite{almahaqeri2026gradient} & CICIoT2023 (46 pož.) & LightGBM &
  dvejetainė & tikslumas 99,53~\%; macro-F1 0,9944;
  \SI{0.72}{\micro\second} &
  Skirtumas nuo XGBoost mažesnis nei tikėtinas matavimo svyravimas \\
\addlinespace
\cite{nassef2026tinyml} & CICIoT2023 & GAT + BiGRU, federuotas
  mokymasis & nenurodyta & F1 0,92; \SIrange{120}{180}{\milli\second}
  Raspberry Pi 4 &
  Užduoties granuliarumas nenurodytas; delsa 2,4--9 kartų viršija
  \SIrange{20}{50}{\milli\second} šliuzo biudžetą \cite{sallam2026gap} \\
\cite{nassef2026tinyml} & Edge-IIoTset & GAT + BiGRU &
  nenurodyta & F1 0,94 &
  Ta pati problema; pramoninis rinkinys, kitas įrenginių profilis \\
\cite{nassef2026tinyml} & RT-IoT2022 & GAT + BiGRU &
  nenurodyta & F1 0,93 &
  Ta pati problema \\
\addlinespace
\cite{mazinani2026constrained} & CICIoT2023 (46,7~mln. įrašų) &
  MLP, TCN, CNN, LSTM, GRU, SimpleRNN & dvejetainė, 8 ir 34 klasės &
  MLP ir TCN --- geriausias tikslumo ir kainos kompromisas;
  INT8 kvantavimas mažina modelį \textgreater\,90~\% &
  Skelbiamas suminis kompromiso balas, o ne metrikos kiekvienam
  modeliui, todėl tiesiogiai palyginti neįmanoma \\
\addlinespace
\cite{meidan2018nbaiot} & N-BaIoT (9 įrenginiai) & gilusis
  autokoderis & dvejetainė, kiekvienam įrenginiui atskirai &
  TPR 100~\%; FPR $0{,}007 \pm 0{,}01$; aptikimo laikas
  $174 \pm \SI{212}{\milli\second}$ &
  TPR 100~\% pasiektas lengviausioje įmanomoje formuluotėje ---
  vienas įrenginys, dvi botnetų šeimos; aptikimo laiko standartinis
  nuokrypis didesnis už vidurkį \\
\addlinespace
\cite{eren2026drift} & BoT-IoT, ToN-IoT, UNSW-NB15 & įvairūs
  klasifikatoriai & kryžminis perkėlimas &
  vidutinis MCC krinta nuo ${\sim}94~\%$ iki ${\sim}29~\%$ &
  Viename rinkinyje gauti skaičiai kitame negalioja \\
\end{xltabular}
\endgroup

% ---------------------------------------------------------------------
%  BUKLE 2026-09-02, 15:10: uzpildyta 10 eiluciu is 5 saltiniu.
%
%  LIEKA 1 EILUTE:
%   houichi2025smartcity -- CICIoT2023, klasikiniai MM metodai.
%     Leidejo puslapis grazina 403, ResearchGate 429. Reikia
%     universiteto prieigos. Uzpildyti: metodas, uzduotis, metrika.
%
%  NEIEJO I LENTELE (sazininga pataisa):
%   imani2025imbalance -- naudoja telekomunikaciju klientu nutekejimo
%     duomenis su dirbtinai nustatytais disbalanso lygiais (15/10/5/1 %),
%     NE IoT idsibrovimu rinkini. Todel tai NE susijes darbas, o
%     metodologinis saltinis -> 2.6 poskyris (balansavimo pasirinkimas).
% ---------------------------------------------------------------------
% ---------------------------------------------------------------------
% 2.8 ATSISAKYTA (2026-09-02): "galiojimo ribos be kryzminio
%   patikrinimo" isbraukta. Priezastis - vieta, ne apimtis: 2 skyrius
%   apzvelgia metodus, o rezultatu galiojimo ribos priklauso
%   metodikai (4 sk.) ir apribojimu aptarimui (6 sk.). Antriniu
%   rinkiniu tiesiog nera - to konstatavimas nereikalauja poskyrio.
%   Perkelta: viena pastraipa i 4 sk. metodika (skaidymo sprendimas).
\subsection{Skyriaus apibendrinimas: kandidatų aibė}
\label{sec:kandidatai}

Apžvelgti aštuoniolika metodų (žr. \ref{tab:metodai} lentelę).
Kandidatų aibę labiausiai siaurina naudojamų duomenų sandara. Galutinę
atranką su kriterijų svoriais atlieka trečias skyrius; čia fiksuojama,
iš ko renkamasi.

\subsubsection{Keturi pasirinkti metodai}

\textbf{Atsitiktinis miškas} imamas kaip atskaitos modelis. Jis
nereikalauja kruopštaus derinimo, atsparus persimokymui, o klasių svoriai
leidžia tvarkyti disbalansą nedubliuojant duomenų. Jis nustato
atskaitos lygį, prie kurio matuojamas sudėtingesnių metodų duodamas
prieaugis.

\textbf{Gradientinio stiprinimo modelis} (XGBoost) yra laukiamas
geriausias prižiūrimas metodas: tame pačiame duomenų rinkinyje jis
pranoko dešimt kitų metodų, o inferencijos trukmė matuojama
mikrosekundėmis \cite{almahaqeri2026gradient}. Kartu jis yra vienintelis
tiesiogiai palyginamas atskaitos taškas literatūroje.

\textbf{Daugiasluoksnis perceptronas} įtraukiamas kaip vienintelė gilaus
mokymosi architektūra, prasminga lentelinei įvesčiai --- ji nedaro
prielaidų nei apie požymių eilę, nei apie kaimynystę, nei apie seką.
Ribotų išteklių įrenginiuose jis kartu su laikine konvoliucija rodo
geriausią tikslumo ir kainos santykį \cite{mazinani2026constrained}.
Jo paskirtis --- patikrinti, ar didesnis modelio sudėtingumas duoda
naudos ten, kur įvestis yra trisdešimt šešios nepriklausomos
statistikos.

\textbf{Autokoderis} yra vienintelis pasirinktas neprižiūrimas metodas.
Jis mokomas tik iš gerybinio srauto, todėl gali aptikti atakas, kurių
pavyzdžių mokymo metu nebuvo --- tai vienintelis būdas patikrinti
darbe keliamą prielaidą apie nematytų grėsmių aptikimą
\cite{meidan2018nbaiot}.

Pagrindiniu užduoties detalumu eksperimentams pasirenkamos
\textbf{aštuonios atakų kategorijos}. Toks pasirinkimas atitinka darbe
naudojamą kategorijų žodyną ir sutampa su literatūroje skelbiamu
makro-F1 įverčiu $0{,}8903$ \cite{almahaqeri2026gradient}, todėl
rezultatus bus su kuo palyginti. Dvejetainė ir trisdešimt keturių klasių
formuluotės atliekamos tik geriausiam modeliui.

\subsubsection{Kodėl atmesti likusieji}

Rekurentiniai tinklai ir dėmesio mechanizmu grįstos
architektūros modeliuoja sekas, tačiau naudojamame rinkinyje nėra nei
laiko žymos, nei eilės tarp įrašų --- kiekvienas įrašas jau yra dešimties
arba šimto paketų lango apibendrinimas, todėl sekos, kurią būtų galima
modeliuoti, tiesiog nėra. Grafų neuroniniam tinklui reikia mazgų ir
briaunų, o rinkinyje nėra nei adresų, nei krypties, todėl grafo
sudaryti neįmanoma. Dėl tos pačios priežasties negalimas ir federuotas
mokymasis: be įrenginio identifikatoriaus nėra pagal ką dalyti duomenų
į klientus, o atsitiktinis padalijimas imituotų vienodą pasiskirstymą ir
tirtų ne tai, kas svarbu. Naivųjį Bajeso klasifikatorių atmeta jo paties
prielaida: požymių nepriklausomumas paneigiamas tiksliomis funkcinėmis
priklausomybėmis, aptiktomis šiame rinkinyje.

Atraminių vektorių mašina ir jos vienos klasės atmaina atmetamos dėl
kvadratinio arba kubinio mokymo sudėtingumo, o artimiausių kaimynų
metodas --- dėl inferencijos trukmės, kuri auga kartu su mokymo aibe ir
netelpa į kraštinio šliuzo laiko biudžetą.

Sprendimų medis, lengvasis gradientinio stiprinimo modelis, izoliacijos
miškas ir vienmatis konvoliucinis tinklas neatrinkti dėl pertekliškumo:
jie atkartotų paradigmas, kurioms jau atstovauja pasirinktieji.

\subsubsection{Palyginimo asimetrija}

Keturi pasirinkti metodai nėra vienarūšiai, ir į tai reikia atsižvelgti
planuojant vertinimą. Trys prižiūrimi metodai klasifikuoja į aštuonias
kategorijas, o autokoderis pateikia tik anomalijos įvertį, t.\,y.
sprendimą „ataka arba ne ataka“. Todėl bendro makro-F1 stulpelio
visiems keturiems sudaryti negalima: prižiūrimi metodai lyginami
daugiaklasėje formuluotėje, o autokoderis --- dvejetainėje ir papildomai
nematytų atakų klasių bandymu. Šeštojo skyriaus suvestinė lentelė dėl to
turės dvi dalis.

\medskip
Trečiajame skyriuje šie keturi metodai vertinami pagal pirmame skyriuje
išvestus reikalavimus, jiems priskiriami svoriai ir pagrindžiamas
galutinis pasirinkimas.
